%%% Local Variables:
%%% mode: latex
%%% Tex-engine: XeLaTeX
%%% TeX-master: t
%%% End:

\documentclass{article}

\usepackage{fancyhdr}
\usepackage{extramarks}
\usepackage{amsmath,amsthm,amssymb,amsthm}
\usepackage{enumitem}
\usepackage{tikz}
\usepackage[plain]{algorithm}
\usepackage{algpseudocode}
\usepackage[utf8]{inputenc}
\usepackage{color,soul}
\usepackage{graphicx}
\graphicspath{ {./} }

\usetikzlibrary{automata,positioning}

%
% Basic Document Settings
%

\topmargin=-0.45in
\evensidemargin=0in
\oddsidemargin=0in
\textwidth=6.5in
\textheight=9.0in
\headsep=0.25in

\linespread{1.1}

\pagestyle{fancy}
%%\lhead{\hmwkAuthorName}
\chead{\hmwkClass\: \hmwkTitle}
\rhead{\firstxmark}
\lfoot{\lastxmark}
\cfoot{\thepage}

\renewcommand\headrulewidth{0.4pt}
\renewcommand\footrulewidth{0.4pt}

\setlength\parindent{0pt}

%
% Create Problem Sections
%

\newcommand{\enterProblemHeader}[1]{
    \nobreak\extramarks{}{Problem \arabic{#1} continued on next page\ldots}\nobreak{}
    \nobreak\extramarks{Problem \arabic{#1} (continued)}{Problem \arabic{#1} continued on next page\ldots}\nobreak{}
}

\newcommand{\exitProblemHeader}[1]{
    \nobreak\extramarks{Problem \arabic{#1} (cont.)}{Problem \arabic{#1} continued on next page\ldots}\nobreak{}
    \stepcounter{#1}
    \nobreak\extramarks{Problem \arabic{#1}}{}\nobreak{}
}

\setcounter{secnumdepth}{0}
\newcounter{partCounter}
\newcounter{homeworkProblemCounter}
\setcounter{homeworkProblemCounter}{1}
\nobreak\extramarks{Problem \arabic{homeworkProblemCounter}}{}\nobreak{}

%
% Homework Problem Environment
%
% This environment takes an optional argument. When given, it will adjust the
% problem counter. This is useful for when the problems given for your
% assignment aren't sequential. See the last 3 problems of this template for an
% example.
%
\newenvironment{homeworkProblem}[1][-1]{
    \ifnum#1>0
        \setcounter{homeworkProblemCounter}{#1}
    \fi
    
    \section{Problem \arabic{homeworkProblemCounter}}
    \setcounter{partCounter}{1}
    \enterProblemHeader{homeworkProblemCounter}
}{
    \exitProblemHeader{homeworkProblemCounter}
}

%
% Homework Details
%   - Title
%   - Due date
%   - Class
%   - Section/Time
%   - Instructor
%   - Author
%

\newcommand{\hmwkTitle}{Assignment \#1}
\newcommand{\hmwkDueDate}{September 19th, 2018}
\newcommand{\hmwkClass}{COMP 524 - Cybersecurity}
\newcommand{\hmwkClassInstructor}{Dr. Michael Soltys}
\newcommand{\hmwkAuthorName}{
  Rihan Stephen Pereira\\
  \texttt{email - rihanstephen.pereira576@myci.csuci.edu, studentID - 002497665}
  \\
  David Bryne\\
  \texttt{email - something@myci.csuci.edu, studentID - 012345}
  \\
  Marco Cabrera\\
  \texttt{email - something@myci.csuci.edu, studentID - 012345}
  \\
  Kaveh Arashvand\\
  \texttt{email - something@myci.csuci.edu, studentID - 012345}
}
%
% Title Page
%

\title{
    \vspace{2in}
    \textmd{\textbf{\hmwkClass:\ \hmwkTitle}}\\
    \normalsize\vspace{0.1in}\small{Due\ on\ \hmwkDueDate}\\
    \vspace{0.1in}\large{\textit{\hmwkClassInstructor\ \hmwkClassTime}}
    \vspace{3in}
}

\author{\hmwkAuthorName}

\renewcommand{\part}[1]{\textbf{\large Part \Alph{partCounter}}\stepcounter{partCounter}\\}

%
% Various Helper Commands
%

% Useful for algorithms
\newcommand{\alg}[1]{\textsc{\bfseries \footnotesize #1}}

% For derivatives
\newcommand{\deriv}[1]{\frac{\mathrm{d}}{\mathrm{d}x} (#1)}

% For partial derivatives
\newcommand{\pderiv}[2]{\frac{\partial}{\partial #1} (#2)}

% Integral dx
\newcommand{\dx}{\mathrm{d}x}

% Alias for the Solution section header
\newcommand{\solution}{\textbf{\large Solution}}

\newcommand{\justif}[2]{&{#1}&\text{#2}}

% Probability commands: Expectation, Variance, Covariance, Bias
\newcommand{\E}{\mathrm{E}}
\newcommand{\Var}{\mathrm{Var}}
\newcommand{\Cov}{\mathrm{Cov}}
\newcommand{\Bias}{\mathrm{Bias}}
\newcommand{\contradiction}{
  \ensuremath{{\Rightarrow\mspace{-2mu}\Leftarrow}}
}

\begin{document}
\maketitle
\pagebreak

% problem 1
\begin{homeworkProblem}
  Decrypt the following four documents. To submit the assignment, write a two page PDF report, which explains, with code snippets, how you managed to decrypt the documents, and what the documents were.
  \begin{itemize}
    \item \textcolor{cyan}{\ul{assignment-1-a.txt}}: Caesar cipher Base64 encoded text
    \item \textcolor{cyan}{\ul{assignment-1-b.txt}}: MAC cipher Base64 encoded text
    \item \textcolor{cyan}{\ul{assignment-1-c.txt}}: Caesar cipher Base64 encoded jpeg
    \item \textcolor{cyan}{\ul{assignment-1-d.txt}}: MAC cipher Base64 encoded jpeg
  \end{itemize}
  \\
  \textbf{Solution:}


The approach I used to solve this problem was a known text attack. By inspecting the other image that was decoded in problem 3 I was able to determine the header and byte offset that was included when the image was taken from an iphone. Then after we base64 encoded that known image I was able to count the exact offset of the header in base64 characters. Once I had this offset decrypting the key was just a matter of comparing the character at the same offset in both the encrypted file and the know file. Once we had a mapping for every character we were able to write a python script that decrypted and decoded the file. 

\includegraphics{solution_1d.jpg}

The tool I used to calculate the byte offset of the image header was 0xd

The key was: 2JzvLTdIARk3nyDg9s0NhxjGf18YQrq/ei6ZM+ObWCKotaclw7FE5BXHV4puUmPS

The python program solution_1d.py will output the decrypted image to a file. You can run this code: python solution_1d.py  
 

\end{homeworkProblem}
\end{document}

%%% Local Variables:
%%% mode: latex
%%% TeX-master: t
%%% End: